\documentclass[10pt,a4paper]{amsart}
\usepackage[margin=19mm]{geometry}
\usepackage{amsmath,amssymb,mathtools}
\usepackage[scr=rsfs,cal=euler]{mathalfa}
\usepackage{xfrac}
\usepackage[unicode,hidelinks,bookmarks=false]{hyperref}
\newcommand{\ie}{i.e.\ }
\newcommand{\cf}{cf.\ }
\newcommand{\resp}{respectively\ }
\newcommand{\Subsection}{\subsection}
\providecommand{\nobreakdash}{\nobreak\hbox{-}}
\setlength{\emergencystretch}{2em}
\multlinegap0pt
\usepackage{float}
\usepackage{amssymb}
\usepackage[shortlabels]{enumitem}
\usepackage[all]{xy}
\usepackage{stackengine}
\newtheorem{prp}{Proposition}
\newtheorem{lmm}{Lemma}
\newtheorem{thm}{Theorem}
\title{Weyl's Law for compact rank one symmetric spaces}
\author{Sai Sriharsha Indukuri, Ritwik Mukherjee}
\date{Source: arXiv:2407.05274v2 (CC BY 4.0)}
\begin{document}
\maketitle

\noindent\emph{Source and licence.} Weyl's Law for compact rank one symmetric spaces. Version arXiv:2407.05274v2 (CC BY 4.0), distributed by arXiv under the Creative Commons Attribution 4.0 International licence, http://creativecommons.org/licenses/by/4.0/, as recorded in the arXiv licence metadata for this exact version and shown on the abstract page https://arxiv.org/abs/2407.05274v2. Full licence notice: \texttt{licenses/arxiv-2407.05274-CC-BY-4.0.txt}.\par\smallskip

\noindent\textit{Editorial scope.} Counted unit: the article's thm Auxiliary\_result\_2 (source0.tex lines 1430--1461). The article's complete original proof of it is delivered with it (source0.tex lines 1463--1473, one \texttt{proof} environment of 11 lines). No mathematics is written, repaired or re-derived: the statement and the proof are the article's own lines, in the article's own notation, and no retained line is substituted or re-flowed.  Retained uncounted as the source-local context the proof needs: lines 1280--1427.  Nothing was omitted from any retained span, so the package carries no omission record.  No retained line contains a citation, so the wrapper carries no bibliography and no bibliography entry.  No retained line contains a figure environment, an image-inclusion command or drawing code, so no image is delivered and the package declares no asset dependency.  Editorial formatting only, in addition: an AMS \texttt{amsart} preamble that restates the article's own declarations and macros used by the retained lines, namely the environments \texttt{prp}, \texttt{lmm}, \texttt{thm} on separate counters, together with the standard packages those lines need.  The retained environments print in this document as Prp 1, Lmm 1, Theorem 1 in order of appearance, whereas the article prints the same items with the numbering of its own full document.  The complete native source of the article as distributed in the arXiv e-print for this exact version is bundled unchanged as \texttt{sources/161919/source0.tex}.
\begin{prp}
\label{psi_k}
Let us define the following function
$\psi_k: \mathbb{N}\longrightarrow \mathbb{R}$, given by   
\begin{align} 
\psi_k(n)&:= \sum_{i=1}^{n} i^{k-1} \nonumber \\ 
          & = 1^{k-1} + 2^{k-1} + \ldots + n^{k-1},   
\end{align} 
where $k \in \mathbb{Z}^{\geq 2}$. Then, 
\begin{align} 
\psi_k(n)& = \frac{n^k}{k} + \frac{1}{2} n^{k-1} + \mathcal{R}_{k-2}(n), \label{nk_error_sharp}
\end{align}
where $\mathcal{R}_{k-2}(n)$ is a polynomial in $n$ of degree at most $k-2$. 
\end{prp}
%\begin{prp} 
%Let us define the following function $\psi: \mathbb{N}\longrightarrow \mathbb{R}$, given by 
%\begin{align} 
%\psi(n)&:= \sum_{i=1}^{n} i^{k-1} \nonumber \\ 
%          & = 1^{k-1} + 2^{k-1} + \ldots + n^{k-1}.   
%\end{align}
%Then $\psi(n)$ is a polynomial of degree $k$, whose coefficient of $n^{k-1}$ is non zero. Here $k \in \mathbb{Z}^{\geq 2}$. 
%Let us now extend $\psi$ on $\mathbb{R}^{\geq 0}$ by declaring it to be a constant on $[n, n+1)$. 
%We will call the extended function $N$. Then,  
%\begin{align}
%N(\lambda) & = \frac{1}{k}[\lambda]^k + \frac{1}{2}[\lambda]^{k-1}+\text{polynomial in $[\lambda]$ of degree $(k-2)$}.
%\end{align}
%\end{prp}
\begin{proof}
Let $n \in \mathbb{N} $ and $k\in \mathbb{Z}^{\geq 2} $. 
%For $k\in \mathbb{Z}^{\geq 2} $,  
%consider \[S_{k}(n):= \sum_{i=1}^{n} i^{k-1} . \]
Note that
\begin{equation}
\label{tel_s}
\begin{aligned}
&(n+1)^{k}-n^{k}= kn^{k-1}+\frac{k(k-1)}{2} n^{k-2}+ \cdots + 1\\
&n^{k}-(n-1)^{k}= k(n-1)^{k-1}+\frac{k(k-1)}{2} (n-1)^{k-2}+ \cdots + 1\\
&\cdots \\
&\cdots \\
&2^{k}-1^{k}= k \cdot 1^{k-1}+ \frac{k(k-1)}{2} \cdot 1^{k-2} + \cdots+ 1.
\end{aligned}
%\notag
\end{equation}
%and hence, we have 
Adding up all the terms on 
the left hand side of 
%and right hand side of 
%of all 
%the individual 
equation 
\eqref{tel_s}, we get a telescopic sum which equals 
$(n+1)^{k}-1^{k}$. 
Equating this with the summation of the right hand 
side of equation \eqref{tel_s}, 
we conclude that     
%\eqref{tel_s}, **. 
%By adding up this telescopic sum, we conclude that  
\begin{equation}
(n+1)^{k}-1^{k} = k\psi_{k}(n)+\frac{k(k-1)}{2} \psi_{k-1}(n)+\cdots+\psi_1(n). 
\label{nk_sum}
%\notag
\end{equation}
In equation \eqref{nk_sum}, we note that $\psi_{1}(n)=n $.
Using the fact that $\psi_2(n) = \frac{n(n+1)}{2}$ 
and equation \eqref{nk_sum}, equation \eqref{nk_error_sharp} 
follows by using induction on $k$. 
%this identity, and using induction on $k$, 
%we conclude that that for all $k\in \mathbb{Z}^{\geq 2}$,
%\[\psi_{k}(n)=\frac{1}{k} n^{k}+\frac{1}{2} n^{k-1}+ \text{polynomial in $n$ of degree 
%at most $k-2$} ,\]
%This proves the desired the result. 
%follows. 
This proves the desired result.  
%of the proposition.  
\end{proof}
Next, we prove the following result. 
\begin{lmm}
\label{sharp_main}
Let $a \in \mathbb{R}_{>0}$, $b\in \mathbb{R}$ and $y \in \mathbb{R}_{\geq 0}$. For $\lambda$ 
sufficiently large and positive, define 
\begin{align*}
M(\lambda)& = \frac{[\sqrt{a\lambda^2+b}-y]^d}{d} + (y+\frac{1}{2})[\sqrt{a\lambda^2+b}-y]^{d-1}.
\end{align*}
Then
\begin{align}
\lim_{\lambda\longrightarrow +\infty} \frac{M(\lambda)}{\lambda^d} & = \frac{a^{\frac{d}{2}}}{d}. 
\label{M_leading}
\end{align}
Furthermore, the error term $\mathcal{E}(\lambda)$ can not be $O(\lambda^{d-1-\varepsilon})$ 
for any positive $\varepsilon$, where 
\begin{align*}
\mathcal{E}(\lambda) &:= M(\lambda)-\frac{a^{\frac{d}{2}}}{d} \lambda^d.
\end{align*}
\end{lmm} 
\noindent \textbf{Proof:} Equation \eqref{M_leading} is immediate using the fact that 
\begin{align} 
A+B-1&\leq [A+B]\leq A+B+1. \label{A_B_box}
\end{align}
An application of \eqref{A_B_box}, binomial theorem 
and the squeeze rule immediately gives us \eqref{M_leading}.  
We now prove the second assertion. We will prove by contradiction. Suppose 
$\mathcal{E}(\lambda)$ was in $O(\lambda^{d-1-\varepsilon})$ for some $\varepsilon>0$. 
Let us define the following discrete subset of $\mathbb{R}$ given by 
\begin{align*}
\Lambda&:= \Big\{ \lambda\in \mathbb{R}_{>0}: a \lambda^2+b>0, \sqrt{a\lambda^2+b}-y \in \mathbb{Z}_{\geq 0}
\Big\}.
\end{align*}
Define $\gamma:\Lambda\longrightarrow \mathbb{Z}_{>0}$, given by 
\begin{align} 
\gamma(\lambda)&:= \sqrt{a\lambda^2+b}-y. \label{gamma_defn_in_terms_of_lambda} 
\end{align}
To lighten the notation, 
we will abbreviate $\gamma(\lambda)$ by $\gamma$; however, it should always 
be understood that $\gamma$ depends on $\lambda$ via equation \eqref{gamma_defn_in_terms_of_lambda}. 
Hence, if the error term was in $O(\lambda^{d-1-\varepsilon})$, then it would imply that there is some constant 
$C>0$ such that for all $\lambda \in \Lambda$, sufficiently large, we have 
\begin{align}
\Big| \Big(y+\frac{1}{2}\Big) \gamma^{d-1} 
+\Big(\frac{\gamma^d}{d}-\frac{a^{\frac{d}{2}}}{d} \lambda^d\Big) \Big| & \leq C \lambda^{d-1-\varepsilon}. \label{lm_gmm_contrd}
\end{align} 
Using the fact that $|A|-|B|\leq |A+B|$
(note that $|A+B-B| \leq |A+B| + |B|$ which implies what we have written), 
we conclude that 
\begin{align}
\Big|\Big(y+\frac{1}{2}\Big) \gamma^{d-1}\Big| 
-\Big| \frac{\gamma^d}{d}-\frac{a^{\frac{d}{2}}}{d} \lambda^d \Big| 
& \leq C \lambda^{d-1-\varepsilon}. \label{lm_gmm_contrd_ag}
\end{align} 
Here we set 
\[A:= \Big(y+\frac{1}{2}\Big) \gamma^{d-1} \qquad \textnormal{and} 
\qquad B:= \Big(\frac{\gamma^d}{d}-\frac{a^{\frac{d}{2}}}{d} \lambda^d\Big). \]
Now we divide both sides of equation \eqref{lm_gmm_contrd_ag} 
and take the limit as $\lambda$ tends to infinity (and $\lambda \in \Lambda$). 
The right hand side will obviously be zero. We first note that 
\begin{align*} 
\lim_{\substack{\lambda \longrightarrow \infty\\ \lambda \in \Lambda}} 
\frac{\gamma^d}{d} - \frac{a^{\frac{d}{2}}}{d} \lambda^d & = 0. 
\end{align*}
This follows by taking a binomial expansion of $\gamma$ and noticing that after subtracting the 
leading term (namely $\frac{a^{\frac{d}{2}}}{d} \lambda^{d}$), we are left with a terms of order $\lambda^{d-2}$ or less. 
Hence, what remains is the first term on the left hand side of equation 
\eqref{lm_gmm_contrd_ag}. 
We note that 
\begin{align*} 
\lim_{\substack{\lambda \longrightarrow \infty\\ \lambda \in \Lambda}} 
\Big(y+\frac{1}{2}\Big) \frac{\gamma^{d-1}}{\lambda^{d-1}} & = \Big(y+ \frac{1}{2}\Big).
\end{align*} 
Since $y$ is non negative, the above limit is positive. This gives us a contradiction. \qed 

\begin{thm}
\label{Auxiliary_result_2}
Let $d \in \mathbb{Z}^{\geq 2}$  
and $\mathsf{Q}_{d-1}(t)$ be a degree $d-1$ polynomial in $t$ that does not have the $t^{d-2}$ term. In other words, 
\begin{align}
\mathsf{Q}_{d-1}(t) & = C t^{d-1} + \varphi(t), \qquad \textnormal{if} \qquad d>2 \qquad \textnormal{and} \nonumber \\ 
\mathsf{Q}_{d-1}(t) & = C t, \qquad \textnormal{if} \qquad d=2, \nonumber 
\end{align}
where $\varphi(t)$ is a polynomial of degree at most $d-3$ and 
$C$ is some non-zero constant. 
%, that is, 
%\[\mathsf{Q}_{d-1}(t)= C_{1} t^{d-1}+\text{polynomial of degree less than or equal to $(d-3)$ }, \ C_{1}\neq 0. \]
%, that does not have a $t^{d_i-1}$ term, i.e.  
%\begin{align}
%\mathsf{Q}_{d_i}(t)&:= C_i t^{d_i} + \psi_i(t),
%\end{align}
%where $C_i$ is a constant and $\psi_i(t)$ is a polynomial of degree at most $d_i-2$.
Let $y$ be a fixed non negative real number and 
$A$ a fixed positive real numbers. 
Define 
$N:\mathbb{R}\longrightarrow \mathbb{R}$ 
as  
%Given a real number $\lambda$, let $N(\lambda)$ be defined as   
\begin{align}
N(\lambda) & := \sum_{\substack{m-y \in \mathbb{Z}^{\geq 0}\\ A m^2 \leq a \lambda^2 + b}}
\mathsf{Q}_{d-1}(m). 
\end{align} 
where $a \in \mathbb{R}^{>0}$ and $b \in \mathbb{R}$. 
Then the error term in the asymptotic expression of $N(\lambda)$ 
(after subtracting off the leading term) can not be improved by $O(\lambda^{d-1-\varepsilon})$ 
for any positive $\varepsilon$. 
\end{thm} 

\begin{proof} Without loss of generality, let us set $C$ to be equal to one.  
We sum up the leading term of $\mathsf{Q}_{d-1}(t)$ 
(namely the degree $t^{d-1}$ terms) using Proposition \ref{psi_k}. 
The first two terms precisely correspond to $M(\lambda)$ as given by 
Lemma \ref{sharp_main} (replace $a$ by $\frac{a}{A}$ and replace $b$ by $\frac{b}{A}$). 
The remaining terms are all of order $\lambda^{d-2}$ or less 
(this is where we use the fact that $\mathsf{Q}_{d-1}(t)$ does not have a $t^{d-2}$ term). 
Hence the remaining terms can not compensate 
for the error term from $M(\lambda)$. We have already seen that the error term for $M(\lambda)$ 
can not be improved. This shows that the error term for $N(\lambda)$ can not be improved.  
\end{proof}

\end{document}
